\chapteropen{06}

\begin{cele}
  \item opisać nawyk czytania jako trzy elementy: wyzwalacz, miejsce i minimalną dawkę, i sprawdzić, który z nich u ciebie jest najsłabszy,
  \item ułożyć plan \glqq jeśli--to\grqq{} na dni, w których nie da się czytać tak, jak zwykle,
  \item przykleić czytanie do czynności, która już się dzieje codziennie, zamiast polegać na pamięci,
  \item używać licznika serii tak, żeby nie zamienił przerwy w porzucenie książki,
  \item wrócić po kilku dniach przerwy bez nadrabiania zaległości.
\end{cele}

Większość postanowień o czytaniu nie upada z powodu braku chęci. Upada, bo nie ma jasno określonego momentu, w którym czytanie się zaczyna, i bo \glqq kiedyś\grqq{} nie ma żadnej miary. Postanowienie \glqq będę czytać więcej\grqq{} nie mówi, co zrobić o 21:40 we wtorek, kiedy zmęczenie wygrywa z ambicją. Nawyk to coś innego: czynność, którą uruchamia konkretna sytuacja, wykonuje się w znanym miejscu i jest na tyle mała, że trudno wymyślić powód, żeby jej nie zrobić.

Ten rozdział nie obiecuje, że nawyk powstanie w określonym czasie ani że będzie przyjemny od pierwszego dnia. Pokazuje mechanikę, która pozwala ją zbudować, i koszty, które te mechanizmy mają. Najbardziej praktyczna część to wersja minimalna: dziesięć minut dziennie. Jej rachunek jest prosty, a jej wynik jest mniej imponujący, niż chciałby ktoś, kto szuka planu na dwieście książek, i właśnie dlatego zwykle działa.

\section{Wyzwalacz, miejsce i dawka}

Każdy nawyk ma trzy części. Pierwsza to wyzwalacz, czyli sygnał, który uruchamia czynność. Może być czasowy (\glqq o 20:30\grqq{}), miejscowy (\glqq gdy siadam w fotelu\grqq{}) albo zdarzeniowy (\glqq gdy zamknę laptopa\grqq{}). Druga to miejsce, w którym czynność się odbywa. Trzecia to dawka: ile dokładnie robisz, zanim uznasz, że dzień jest zaliczony.

Najczęstszy błąd polega na tym, że dawka jest ustawiona zbyt wysoko. Czytelnik postanawia czytać godzinę dziennie, pierwszego dnia czyta czterdzieści minut, drugiego czyta dwadzieścia, a trzeciego nie czyta wcale. Po tygodniu nie ma poczucia, że cokolwiek się zbudowało, bo każdy dzień kończy się brakiem. Dawka powinna być tak mała, żeby wykonanie było możliwe nawet w gorszy dzień. To nie jest rezygnacja z ambicji. To zmiana jednostki, w której ambicja jest mierzona.

Sprawdź swój nawyk na trzech pytaniach. Czy wiesz, co dokładnie uruchamia czytanie? Czy wiesz, gdzie to się dzieje? Czy wiesz, ile musisz zrobić, żeby dzień się liczył? Jeśli na którekolwiek odpowiedź brzmi \glqq zależy\grqq{}, właśnie tam jest słaby punkt.

\begin{przyklad}[Trzy nawyki, jeden tydzień]
  Anna ma postanowienie \glqq czytać wieczorem\grqq{}. Wyzwalacz: nieokreślony wieczór. Miejsce: sypialnia albo kanapa, zależnie od dnia. Dawka: \glqq tyle, ile się uda\grqq{}. Po tygodniu odbyło się czytanie w dwóch dniach, a w pozostałych pięciu nie ma wiadomo, co się stało, bo nic nie było do sprawdzenia.

  Zmieniamy tylko wyzwalacz i dawkę. Wyzwalacz: zaraz po położeniu telefonu na ładowarkę w kuchni. Miejsce: fotel przy oknie. Dawka: dziesięć minut, mierzone zegarem, a nie stronami. Po tygodniu przeczytała codziennie, a dwa razy przeczytała dłużej, niż wymagała dawka. Zmieniło się nie tyle jej zaangażowanie, ile to, że wiedziała, kiedy zaczyna i kiedy kończy.
\end{przyklad}

Dawka mierzona czasem jest uczciwsza niż mierzona stronami. Strony mają różną gęstość: rozdział o algebrze liniowej zajmuje więcej uwagi niż wspomnienia podróżnika, a dwie strony pierwszej mogą być warte dwóch razy więcej niż dwie strony drugiej. Zegar nie zależy od gęstości tekstu, więc daje ci wynik, który da się porównać z następnym dniem. Jeśli wolisz strony, licz je jako dawkę zastępczą, ale zmierz najpierw, ile słów mieści się na twojej stronie. Typowa książka ma \val{typical.words} słów, a liczba stron zależy od wydania, czcionki i marginesów.

\section{Plan \glqq jeśli--to\grqq{} na dni awaryjne}

Najlepiej działające plany nie zakładają, że każdy dzień będzie dobry. Zakładają, że wiesz, co zrobisz w dniu gorszym, zanim on nadejdzie. Taki plan ma formę zdania: \glqq jeśli zdarzy się X, to zrobię Y\grqq{}. Badania nad planowaniem wykonania (ang. implementation intentions) opisują ten mechanizm, a jego skuteczność zależy od zadania i od tego, jak konkretny jest sygnał. \verify{skala i powtarzalność efektu w badaniach nad planami \glqq jeśli--to\grqq{} różni się między zadaniami i nie należy przenosić jej na czytanie bez sprawdzenia}.

Ważne jest to, co mieści się w planie awaryjnym. Dobry plan nie wymaga pełnej dawki. Wymaga minimum, które zachowuje ciągłość. Oto kilka przykładów, które da się od razu dopisać do własnej listy:

\begin{itemize}
  \item Jeśli wracam z pracy tak późno, że nie mam dziesięciu minut, to czytam jedną stronę w kolejce albo w windzie.
  \item Jeśli jestem chory i nie mogę usiąść, to słucham rozdziału audio albo przeglądam spis treści, żeby nie zerwać kontaktu z książką.
  \item Jeśli w domu jest gość, to czytam dziesięć minut przed jego przyjściem albo po wyjściu, a nie w trakcie rozmowy.
  \item Jeśli jestem w podróży bez książki, to czytam to, co mam na telefonie, nawet jeśli to nie jest moja lektura.
\end{itemize}

Zwróć uwagę, że w każdym z tych przypadków jest nie tylko alternatywa, ale też granica. Plan nie mówi \glqq jeśli jest źle, to nie czytam\grqq{}. Mówi \glqq jeśli jest źle, to robię mniej, ale czytam\grqq{}. Różnica między tymi dwoma zdaniami jest większa, niż wygląda: pierwsze zamyka ciągłość, drugie ją chroni.

\begin{zasada}[Plan awaryjny ma być mniejszy od planu zwykłego]
  Jeśli plan awaryjny jest równie wymagający jak zwykły, w gorszy dzień nie będzie wykonany i zostanie porzucony. Jeśli jest o wiele mniejszy, przestaje być awaryjny i staje się zwykłym nawykiem. Dobry plan awaryjny to ta sama czynność w czterech razy mniejszej dawce.
\end{zasada}

\section{Doklejanie do istniejącego nawyku i książka w zasięgu}

Nowy nawyk łatwiej zbudować, przyczepiając go do czynności, która już się odbywa codziennie. Kawa rano, mycie zębów wieczorem, odpalenie komputera po wejściu do biura: to są wyzwalacze, których nie trzeba pamiętać, bo ciało wykonuje je bez decyzji. Czytanie dołączone do takiej czynności przestaje wymagać postanowienia za każdym razem.

Schemat jest prosty: \glqq po zaparzeniu kawy otwieram książkę na zaznaczonej stronie i czytam dziesięć minut\grqq{}. Druga część zdania jest najważniejsza, bo daje miejsce i dawkę. Pierwsza część daje wyzwalacz. Sama kawa nie wystarcza, jeśli nie ma określonego miejsca i określonej książki. Dlatego dobry dopisek brzmi: \glqq po zaparzeniu kawy siadam w fotelu przy oknie i otwieram książkę tam, gdzie zostawiłem zakładkę\grqq{}.

Osobnym problemem jest dostępność książki. Książka, która leży w innym pokoju, wymaga decyzji: wstać, poszukać, przynieść. Każda taka decyzja to szansa na rezygnację. Dlatego książka powinna być tam, gdzie zwykle jest wyzwalacz: na stoliku przy ekspresie, na poduszce, w torbie, która jedzie z tobą do pracy. Jeśli masz czytnik albo aplikację, zostaw je na ekranie głównym, a nie w folderze. Jeśli masz dwie książki, na wierzchu leży ta, którą czytasz teraz, a nie ta, którą chciałbyś czytać.

Mycie zębów to jedna z najpopularniejszych kotwic w tego typu poradach, ale wybór kotwicy jest mniej ważny niż to, czy ją naprawdę da się zauważyć. Wybierz czynność, którą wykonujesz codziennie, o podobnej porze, bez zależności od pogody czy dnia tygodnia. Jeśli taka czynność nie istnieje w twoim dniu, nie wymyślaj jej na siłę. Wtedy lepszy jest wyzwalacz czasowy, który sam ustawisz.

\begin{falle}[Kotwica, która nie jest codzienna]
  Czytanie \glqq po przyjściu do domu\grqq{} brzmi jak dobry wyzwalacz, dopóki nie zauważysz, że w połowie tygodnia wracasz do domu o różnych porach albo wcale. Kotwica musi być czymś, co dzieje się w każdym dniu, który chcesz liczyć. Jeśli wyzwalacz zdarza się w sześciu dniach z siedmiu, zaplanuj także siódmy: co robisz w niedzielę, kiedy nie ma kawy w pracy?
\end{falle}

\section{Licznik serii i jego pułapka}

Licznik serii, czyli liczba dni z rzędu, w których czytałeś, jest jednym z najczęściej wybieranych narzędzi. Działa, bo daje widoczny wynik i sprawia, że przerwa staje się zauważalna. Ma jednak wadę, której trzeba się przyjrzeć uczciwie. Im dłuższa seria, tym większy koszt jej przerwania, i tym większa pokusa, żeby nie przerywać jej po jednym gorszym dniu, nawet kosztem zdrowego odpoczynku.

Pułapka wygląda zwykle tak. Seria wynosi czterdzieści dni. Pierwszego wieczoru, w którym nie masz siły, zaczynasz rozważać, czy nie przeczytać jednej strony tylko po to, żeby seria trwała. Ta jedna strona wykonuje dokładnie tę samą funkcję co dziesięć minut: utrzymuje ciągłość. Ale jeśli drugiego wieczoru sytuacja się powtórzy, pytanie zmienia się w \glqq czy warto zerwać serię, która jest już tak długa\grqq{}. Licznik zaczyna wtedy sterować decyzją, a nie raportować jej.

Licznik ma sens, jeśli spełnia dwa warunki. Po pierwsze, liczy dni, w których wykonałeś dawkę minimalną, a nie dni, w których czytałeś dużo. Po drugie, przerwa w serii nie oznacza porażki nawyku, tylko zdarzenie, które zapisujesz i na które reagujesz planem awaryjnym. Jeśli nie potrafisz tego zrobić, lepiej nie prowadzić licznika wcale.

\begin{przyklad}[Seria, która nie jest serią]
  Marek prowadzi licznik od miesiąca. Ma czterdzieści dni z rzędu i zapisuje w nich wszystko, co przeczytał, również bardzo krótkie sesje. Jednego wieczoru ma gorączkę. Przeczytanie dziesięciu minut jest fizycznie niemożliwe, więc zamiast tego przeczytał jedną stronę przez łzy i zapisał to jako dzień zaliczony.

  Licznik pokazuje czterdzieści jeden dni, ale Marek wie, że jeden z nich nie był dniem czytania w sensie, który chce śledzić. Zmienia zasadę: dzień liczy się, gdy jest dawka minimalna albo wpis w planie awaryjnym, a seria przerwana chorobą nie jest kasowana, tylko oznaczona. Licznik przestał być narzędziem do karania siebie, a stał się zapisem tego, co faktycznie się działo.
\end{przyklad}

\begin{uwaga}[Licznik nie jest wskaźnikiem efektu]
  Liczba dni z rzędu mierzy regularność wykonania, a nie to, ile książek przeczytasz ani ile zrozumiesz. Długa seria z dziesięciominutowymi sesjami może oznaczać mniej przeczytanych stron niż krótsza seria z dłuższymi sesjami. Traktuj licznik jako dowód, że nawyk działa, a nie jako miarę tego, że działa dobrze.
\end{uwaga}

\section{Powrót po przerwie bez nadrabiania}

Każdy, kto czytał regularnie, w końcu przerywa. Może to być tydzień urlopu, choroba, wyjazd, awaria, okres, w którym cały czas zajmuje coś innego. Problem nie leży w przerwie, tylko w tym, co robisz zaraz po niej. Najczęstsza reakcja to nadrabianie: \glqq straciłem cztery dni, więc przeczytam dwa razy tyle\grqq{}. Ta logika brzmi sensownie, ale dokłada obciążenie w chwili, w której już jest trudno wrócić.

Powrót wymaga zasady, którą łatwo zapamiętać. Zacznij od dawki minimalnej, a nie od miejsca, w którym przerwałeś z ambicją. Nie czytaj brakujących fragmentów, jeśli nie są potrzebne do zrozumienia dalszego tekstu. Nie licz zaległości w minutach ani w stronach. Zapisz, że przerwa była, i wróć do wyzwalacza.

\begin{zasada}[Po przerwie wracasz do dawki minimalnej]
  Liczba przegapionych dni nie zmienia dzisiejszej dawki. Jeśli przerwa wyniosła trzy dni, dzisiaj czytasz dziesięć minut, a nie trzydzieści. Nadrabianie to dług, który spłacasz kosztem następnych dni, a to właśnie on zwykle prowadzi do kolejnej przerwy.
\end{zasada}

\begin{przyklad}[Trzy dni urlopu]
  Iza nie czytała przez trzy dni, bo była na urlopie bez książki. Pierwszego dnia po powrocie mogłaby przeczytać trzydzieści minut, żeby nadrobić. Zamiast tego czyta dziesięć minut, tak jak zwykle, i pomija te fragmenty, które dotyczą jedynie streszczenia tego, co działo się w tym czasie. Trzeciego dnia wraca do zwykłego tempa.

  Rachunek: trzy dni po dziesięć minut to trzydzieści minut, które przepadły. Nadrobienie wymagałoby dodatkowych trzydziestu minut w jednym dniu, czyli w sumie czterdziestu, a to jest dokładnie ten rodzaj obciążenia, który zapowiada kolejne porzucenie. Trzydzieści minut straty nie jest wielką stratą. Dwa dni przeciążenia mogą nią być.
\end{przyklad}

Jest też jedna sytuacja, w której nadrabianie ma sens: gdy książka ma wątek, który wymaga ciągłości, i przerwa rozbiła twoje rozumienie. Wtedy nie czytasz więcej, tylko przypominasz sobie ostatni rozdział, zanim wrócisz do bieżącego. To jest jedna lektura, a nie nadrabianie. Rozróżnienie jest proste: nadrabiasz, gdy chcesz zrekompensować czas; przypominasz, gdy chcesz odzyskać sens.

\section{Nagrody i czas budowania nawyku}

Nagrody są wygodne i ludzie często je stosują. Herbata po sesji, odcinek ulubionego serialu po tygodniu, drobny zakup za serię. Ich zaletą jest to, że dają natychmiastową przyjemność za czynność, która sama w sobie może wymagać wysiłku. Ich wadą jest to, że mogą zmienić motywację: czytanie zaczyna być środkiem do nagrody, a nie tym, po co czytasz. \verify{wpływ zewnętrznych nagród na motywację wewnętrzną jest przedmiotem badań z wynikami zróżnicowanymi i zależnymi od kontekstu; nie należy traktować tego jako ustalonego prawa}.

Praktyczna zasada jest więc ostrożna. Nagroda powinna być mała, od razu po czynności i niezależna od tego, ile przeczytałeś. Jeśli wymyślasz nagrodę za każde przeczytane pięćset stron, zachęcasz do szybkiego czytania, a nie do uważnego. Jeśli nagroda ma być tylko za samo wykonanie dawki, łatwiej ją utrzymać przez miesiące.

Pytanie o czas, w którym nawyk się utrwala, jest równie częste i równie łatwo źle zadane. Często powtarza się liczbę dni, która ma być potrzebna, aby czynność stała się automatyczna. Ta liczba funkcjonuje w obiegu jako pewna, choć jej źródło jest kiepskie. \verify{popularna liczba 21 dni jako czas tworzenia nawyku nie wynika z badania, które można uznać za solidne; prawdziwy czas zależy od czynności, osoby i kontekstu, a rozrzut w badaniach jest bardzo duży}. Dla celów tej książki wystarczy wniosek praktyczny: nie zakładaj, że nawyk jest zbudowany po określonej liczbie dni. Sprawdzaj, czy czynność w ogóle się odbywa, i zwracaj uwagę na to, czy przestaje wymagać decyzji. Ten drugi sygnał jest lepszy niż kalendarz.

\begin{uwaga}[Czego nie obiecujemy]
  Ten rozdział nie obiecuje, że po miesiącu czytanie stanie się przyjemnością ani że po trzech miesiącach nie będziesz musiał się przemóc. Część osób tak ma, część nie, a większość czyta w pewnych dniach z trudem także po długim czasie. Miarą sukcesu jest to, czy czytanie się odbywa, a nie to, czy jest łatwe.
\end{uwaga}

\section{Wersja minimalna: dziesięć minut dziennie}

Najprostsza wersja tego rozdziału to jedno zobowiązanie: dziesięć minut dziennie, w tym samym miejscu, przy tej samej kotwicy. Rachunek jest jasny i nie wymaga żadnej wiary w cokolwiek. Dziesięć minut to 600 sekund, a przy tempie 250 słów na minutę daje to 2500 słów dziennie. W ciągu roku, przy założeniu, że czytasz każdego dnia, daje to:

\begin{przyklad}[Rok z dziesięcioma minutami]
  Dziesięć minut dziennie, 365 dni, 250 słów na minutę, typowa książka o długości \val{typical.words} słów. Wynik: \val{books.10.250} książki rocznie.

  Przy tempie 200 słów na minutę ten sam rachunek daje \val{books.10.200} książki. Zmieniło się tylko tempo, a nie czas ani długość książek.
\end{przyklad}

Wynik \val{books.10.250} książki rocznie wygląda skromnie, i to jest dobra wiadomość. Nie trzeba być wyjątkowo zdyscyplinowanym, żeby go osiągnąć. Wystarczy nie przerywać nawyku na długo i nie zamieniać go w coś, co zabiera godzinę. Jeśli po roku okaże się, że czytasz więcej niż ta liczba, to znaczy, że dawka okazała się zbyt mała i można ją zwiększyć. Zwiększaj ją dopiero wtedy, gdy wykonanie jest pewne, a nie wtedy, gdy masz gorszy tydzień.

Dziesięć minut dziennie ma też jedną przewagę, której nie widać w tabeli: jest to wersja, którą zapamiętujesz. Kiedy ktoś zapyta, co robisz z wieczorami, odpowiedź \glqq czytam dziesięć minut po kawie\grqq{} jest konkretna i sprawdzalna. Konkretna odpowiedź zmniejsza też pokusę, żeby udawać przed sobą, że czytasz więcej.

Na koniec jedna praktyczna myśl. Książka, która leży na stoliku przy ekspresie, nie zmieni twojego życia sama. Ale dziesięć minut, które robisz codziennie w tym samym miejscu, po tej samej czynności, zwykle wystarcza, żeby do końca roku przeczytać więcej książek, niż przeczytałbyś bez nich, i żeby nie musieć o tym decydować co wieczór.

\begin{cwiczenia}
  \exitem{Zapisz swój nawyk czytania jako jedno zdanie z wyzwalaczem, miejscem i dawką. Które z trzech pól jest najmniej konkretne?}{Wzorcowo: \glqq po zaparzeniu kawy siadam w fotelu i czytam dziesięć minut\grqq{}. Najczęściej najsłabsze jest pole dawki (\glqq tyle, ile się uda\grqq{}) albo wyzwalacza (\glqq wieczorem\grqq{}).}
  \exitem{Napisz trzy plany \glqq jeśli--to\grqq{} na dni awaryjne. Każdy musi mieć mniejszą dawkę niż zwykły dzień.}{Przykład: \glqq jeśli nie mam dziesięciu minut, to czytam jedną stronę w kolejce\grqq{}. Dawka awaryjna jest czterokrotnie mniejsza od zwykłej, więc plan działa także w gorszy dzień.}
  \exitem{Masz zwykłą dawkę dziesięć minut i przerwę trzech dni. Ile minut czytasz pierwszego dnia po powrocie i dlaczego?}{Dziesięć minut, czyli dawka minimalna. Nadrabianie dałoby trzydzieści minut tego dnia, a przy zwykłej dawce strata to tylko trzydzieści minut, a nie powód do przeciążenia.}
  \exitem{Przy liczniku serii czterdziestu dni pojawia się gorszy wieczór. Co sprawdzasz, zanim zdecydujesz, czy przerwać serię?}{Czy wykonałeś dawkę minimalną albo plan awaryjny. Jeśli tak, seria trwa. Jeśli nie, przerwa jest zapisana i reagujesz planem awaryjnym na następny dzień, bez karania się.}
  \exitem{Podaj czynność, która już odbywa się codziennie i może być kotwicą dla czytania. Dlaczego nie jest to wyzwalacz \glqq po powrocie do domu\grqq{}?}{Przykład: mycie zębów wieczorem albo zaparzenie kawy rano. Powrót do domu bywa o różnych porach, więc nie zachodzi w każdym dniu, który liczysz.}
  \exitem{Przy tempie 200 słów na minutę i dziesięciu minutach dziennie, ile książek rocznie przeczytasz, zakładając \val{typical.words} słów na książkę?}{$10 \cdot 365 \cdot 200 / 80\,000 \approx 9{,}1$ książki, czyli \val{books.10.200}.}
  \exitem{Dlaczego nagroda za przeczytanie pięciuset stron może zaszkodzić, a nagroda za samo wykonanie dawki zwykle nie?}{Pierwsza nagradza tempo i liczbę stron, co zachęca do szybkiego czytania; druga nagradza regularność, czyli to, co chcesz utrwalić, a nie to, ile się przeczyta.}
\end{cwiczenia}

\zadanie{Zapisz swój nawyk czytania w formie jednego zdania: wyzwalacz, miejsce, dawka, a w osobnej linijce plan awaryjny na dwa najbardziej prawdopodobne gorsze dni. Przez dwa tygodnie wykonuj dawkę minimalną po kotwicy, którą wybrałeś, i nie zwiększaj jej. Miara sukcesu: na koniec drugiego tygodnia prowadzisz zapis, w którym każdy dzień ma jedną z trzech wartości: dawka wykonana, plan awaryjny albo przerwa. Brak wpisu jest przerwą. Jeśli przerw jest więcej niż dwie, zmień kotwicę albo miejsce, a nie dawkę.}
